
\chapter{带局部系数的奇异链}
\label{app:a2}
\liftIndex{I366}

本附录回顾带局部系统系数的奇异链。这里没有新内容；目的只是固定本书所用
的记号与约定。

为满足本书的需要，只讨论空间为流形的情形。设 $M$ 是 $n$ 维光滑流形。
（在应用中，$M$ 或者是拉格朗日子流形 $L$，或者是它们的交
$L_1\cap L_0$ 的一个连通分支；后一情形假设 $L_1,L_0$ 干净相交。）
为简单起见，假设 $M$ 连通。固定基点 $p_0\in M$。设 $R$ 是交换代数，
$V$ 是自由 $R$-模。给定表示
\[
 \rho\colon\pi_1(M,p_0)\longrightarrow GL(V),
\]
可以在 $M$ 上关联一个平坦 $V$-丛 $\Theta(\rho)$；它也称为与 $\rho$
关联的局部系统。

设 $\Delta^k$ 是标准 $k$-单形，并在 $\Delta^k$ 中选定基点 $o$。

\begin{definition}{A2.1}\label{def:a2-1}
若 $f\colon\Delta^k\to M$ 是奇异单形，并且
\[
 \xi\in\left.\Theta(\rho)\right|_{f(o)},
\]
就称 $\sigma=(f,\xi)$ 为一个\emph{带 $\Theta(\rho)$ 系数的奇异
$k$-单形}。
\end{definition}

$\xi$ 可延拓为平坦丛 $f^*\Theta(\rho)\to\Delta^k$ 的平行截面；下文仍以
$\xi$ 表示这一平行截面。对 $a\in R$ 和 $\sigma=(f,\xi)$，把
$a\cdot\sigma$ 与 $(f,a\cdot\xi)$ 视为相同。以 $S_k(M;\rho,R)$ 表示
所有带 $\Theta(\rho)$ 系数的奇异 $k$-链所成集合；其中的元素写成
\[
 c=\sum_i a_i\cdot\sigma_i,
\]
这里 $a_i\in R$，而 $\sigma_i$ 是带 $\Theta(\rho)$ 系数的奇异
$k$-单形。

设 $\partial_j\Delta^k$ 是 $k$-单形 $\Delta^k$ 的第 $j$ 个面。把
$f\colon\Delta^k\to M$ 限制到 $\partial_j\Delta^k$，得到
\[
 \partial_jf\colon\Delta^{k-1}\longrightarrow M.
\]
再利用从 $\Delta^k$ 的基点到
$\Delta^{k-1}\cong\partial_j\Delta^k$ 的基点的道路，以及
$\Theta(\rho)$ 的平坦结构，得到 $\partial_j\xi$。换言之，把
$\Delta^k$ 上的平坦截面 $\xi$ 限制到
$\partial_j\Delta^k\cong\Delta^{k-1}$，就得到 $\partial_j\xi$。令
\[
 \partial_j(f,\xi)=(\partial_jf,\partial_j\xi).
\]
于是定义边界算子
\[
 \partial=\sum_{j=0}^k(-1)^j\partial_j
 \colon S_k(M;\rho,R)\longrightarrow S_{k-1}(M;\rho,R).
\]
于是 $(S_*(M;\rho,R),\partial)$ 是链复形。所得同调群
$H_*(M;\rho,R)$ 称为\emph{带局部系统 $\Theta(\rho)$ 系数的同调}。

一个带 $\Theta(\rho)$ 值的 $k$-上链 $c$，给每个奇异 $k$-单形
$f\colon\Delta^k\to M$ 指定 $f^*\Theta(\rho)$ 的一个平行截面
$c(f)$。以 $(S^*(M;\rho,R),\delta)$ 表示带 $\Theta(\rho)$ 值的上链
复形，以 $H^*(M;\rho,R)$ 表示所得上同调群。

给定两个表示
\[
 \rho\colon\pi_1(M,p_0)\longrightarrow GL(V),\qquad
 \rho'\colon\pi_1(M,p_0)\longrightarrow GL(V'),
\]
张量积表示
\[
 \rho\otimes\rho'\colon\pi_1(M,p_0)
 \longrightarrow GL(V\otimes V')
\]
由
\[
 (\rho\otimes\rho')(\gamma)=\rho(\gamma)\otimes\rho'(\gamma)
\]
定义。令
\[
 \rho_0\colon\pi_1(M,p_0)\longrightarrow\{\pm1\}
\]
为一维表示，其中 $\rho_0(\gamma)=-1$ 当且仅当切丛沿 $\gamma$ 的和乐
反转定向。把相应的局部系统 $\Theta(\rho_0)$ 记作 $\Theta_M$，并称为
流形 $M$ 的定向丛。由于定向丛的结构群是 $\{\pm1\}$，可以在其上定义
范数。

令 $\rho^\dagger=\rho\otimes\rho_0$。为简单起见，假设 $M$ 是无边界的
$n$ 维紧流形。于是有：

\begin{theorem}{A2.2}\label{thm:a2-2}
设 $\Theta(\rho)$ 是 $M$ 上的局部系统，则
\[
 H^\ell(M;\rho,R)\cong H_{n-\ell}(M;\rho^\dagger,R).
\]
\end{theorem}

\begin{example}{A2.3}\label{ex:a2-3}
设 $P$ 是带边界或角的 $k$ 维流形，$f\colon P\to M$ 是光滑映射。假设
有平坦丛同构
\[
 \phi\colon f^*\Theta(\rho)\cong\Theta_P.
\]
取 $P$ 的一个三角剖分 $\{\Delta_a^k\}$，并以
\[
 \iota_a\colon\Delta^k\longrightarrow\Delta_a^k\subset P
\]
表示相应的奇异单形。令
\[
 c=\sum_a\sigma_a
\]
为 $P$ 的一个带 $\Theta_P$ 系数的基本链。这里
$\sigma_a=(\iota_a,\xi_a)$ 是带 $\Theta_P$ 系数的奇异单形，而
$\xi_a$ 是 $\iota_a^*\Theta_P$ 的单位长度平坦截面。注意，$\phi$ 诱导
同构
\[
 (f\circ\iota_a)^*\Theta(\rho)\cong\iota_a^*\Theta_P.
\]
因此，可把 $\xi_a$ 看作 $(f\circ\iota_a)^*\Theta(\rho)$ 的平坦截面。%
\translatornote{源书在求和项中写 $f\circ\iota$，但本段只定义了逐单形
映射 $\iota_a$，且同一项的截面带下标 $a$；译文补回 $\iota_a$。见疑误
A-0717。}
定义
\[
 f_*(c)=\sum_a(f\circ\iota_a,\xi_a)\in S_k(M;\rho,R).
\]

这一构造还可推广到如下情形：$f\colon P\to M$ 配有同构
\[
 \phi\colon f^*\Theta(\rho)\cong\Theta(\rho'),
\]
其中 $\rho'$ 是 $\pi_1(P)$ 到 $GL(V)$ 的表示。
\end{example}

\begin{example}{A2.4}\label{ex:a2-4}
把例~\ref{ex:a2-3} 的情形推广如下。设 $P$ 是带惠特尼分层的空间，
$f\colon P\to M$ 是光滑映射。此时可以选取 $P$ 的光滑三角剖分，并按
例~\ref{ex:a2-3} 的方式构造。若
\[
 f^*\Theta(\rho)\otimes\Theta_{P_{\mathrm{reg}}}
\]
有一个平凡化，就得到 $S_k(M;\rho,R)$ 中的元素。这里
$P_{\mathrm{reg}}$ 是该分层的顶维层。
\end{example}

\begin{example}{A2.5}\label{ex:a2-5}
设 $X$ 是带仓西结构空间，$f$ 是从 $X$ 到 $M$ 的光滑强连续映射。

设 $R$ 是含 $\mathbb Q$ 的交换环，$\Theta(\rho)$ 是 $M$ 上的
$R$-局部系统。假设 $X$ 有切丛。令 $\mathfrak s_\epsilon$ 如定理
\ref{thm:a1-23}，并选取它，使
$\bigl(\mathfrak s_\epsilon^{-1}(0)\bigr)_{\mathrm{set}}$ 有一个三角剖分
（引理~\ref{lem:a1-26}）。在 $V_p$ 上定义实线丛 $\Theta_X$：
\begin{equation}
 \Theta_X=\Lambda^{\mathrm{top}}E_p^*
 \otimes\Lambda^{\mathrm{top}}TV_p.
 \tag{A2.18}\label{eq:a2-18}
\end{equation}
右端记号同定义~\ref{def:a1-17}。固定
$f^*\Theta(\rho)\otimes\Theta_X$ 的一个平凡化 $\xi$。定义
\[
 f_*\bigl(X^{\mathfrak s_\epsilon},\xi\bigr)
 =\sum_a\operatorname{mul}_{\Delta_a^d}(f_a,\xi_a)
 \in S_k(M;\rho,R),
\]
其中 $(f_a,\xi_a)$ 如例~\ref{ex:a2-3}，而
$\operatorname{mul}_{\Delta_a^d}$ 如定义~\ref{def:a1-27}；由这些重数
构成虚基本链的方式见定义~\ref{def:a1-28}。
\translatornote{源书在式 (A2.18) 中把切空间 $TV_p$ 印成 $V_p$；但
$\Lambda^{\mathrm{top}}$ 只能作用于向量空间／向量丛，且紧接着明确回指
定义 A1.17 的 $\Lambda^{\mathrm{top}}E_p^*\otimes
\Lambda^{\mathrm{top}}TV_p$，故译文补回 $T$。末式的
$S_k(M;\rho,R)$ 又漏排 $\rho,R$ 间的逗号，译文依本附录统一记号补回。
源书另称重数 $\operatorname{mul}_{\Delta_a^d}$ 见定义 A1.28；该重数
实际定义于 A1.27，A1.28 定义的是相应虚基本链，译文把两项均列明。
见疑误 A-0718--A-0720。}
\end{example}
